To ensure precise spatial correspondence between MRI and dual-modality PET images, we implemented a standardized preprocessing pipeline. T1-weighted MRI scans were first skull-stripped using FreeSurfer~\cite{Freesurfer} and subjected to histogram matching for contrast normalization. Each PET image was then rigidly registered to its corresponding MRI. Both modalities were jointly affinely normalized to the MNI152 standard space using SPM~\cite{spm,friston1995spatial}, with identical transformation matrices applied to preserve cross-modality alignment. Following spatial normalization, all volumes were cropped to a common field of view of $160 \times 192 \times 160$ voxels ($1\,\text{mm}^3$ isotropic resolution) and intensity-normalized to the $[0, 1]$ range via min–max scaling. For quantitative regional analysis, T1-weighted MRI scans were parcellated into 98 anatomical regions using SynthSeg~\cite{synthseg}, and the resulting ROI masks were applied to the spatially aligned PET images. ADNI provided preprocessed PET images upstream of this pipeline. Spatial registration and intensity normalization do not establish comprehensive scanner-specific harmonization; acquisition and upstream processing details could not be verified for every actual model input.
